\documentstyle[% 


righttag% 


,11pt% 


,amssymb% 


,russian%               remove in Eng. 


,izvru%                 Set izven% in Eng. 


% ,izvru-ks             for brief communications 


]{amsart} 


 


%       Math definitions 


 


%\newcommand{\ov}{\overline} 


%\newcommand{\wt}{\widetilde} 


\newcommand{\re}{\operatorname{Re}} 


\newcommand{\im}{\operatorname{Im}} 


\newcommand{\mes}{\operatorname{mes}} 


\newcommand{\const}{\operatorname{const}} 


\newcommand{\intO}{\iint\limits_{\Omega}} 


\newcommand{\tts}[1]{} 


 


\theoremstyle{definition} 


\theorembodyfont{\rm} 


\newtheorem{defnnonum}{\hskip\parindent Определение} 


\renewcommand{\thedefnnonum}{} 


 


 


%\hoffset=-15mm%                  For Eng. 


  


\setcounter{page}{53} 


\god{1998} 


\nomer{7~(434)} %                        Remove in Eng. 


%\nomer{Vol.\,42, No.\,7}%               Set in Eng. 


%\str{\pageref{begin}--\pageref{end}}%  Set in Eng. 


\udk{539.3} 


\author{\it И.Г.\,Терегулов, С.Н.\,Тимергалиев} 


% NB:   ^^ remove in Eng. 


\title{О разрешимости одной геометрически\\ 


нелинейной задачи теории пологих оболочек} 


 


\thanks{Работа выполнена при финансовой поддержке Российского фонда 


фундаментальных исследований, проект \No\,96-01-00518, и Академии наук 


Татарстана.} 


\date{} 


% \pagestyle{myheadings}%         Set in Eng. 


 


\pagestyle{plain} %              Remove in Eng. 


\begin{document} 


 


% \thispagestyle{plain}%          Set in Eng. 


 


\maketitle 


\label{begin} 


 


Данная работа посвящена исследованию разрешимости задачи геометрически 


нелинейной теории пологих оболочек, заключающейся в определении 


напряженно-деформированного состояния свободных оболочек, не 


подчиненных никаким геометрическим граничным условиям. Необходимость 


изучения задач для таких оболочек была отмечена И.И.\,Воровичем в 


\cite{1}. Для исследования задач предлагается метод, основанный на 


решении задачи в деформациях \cite{2}. 


\medskip 


 


{\bf 1.} В этом пункте будут выведены формулы для вектора перемещения


через компоненты деформаций. Для этого воспользуемся соотношениями для 


компонент конечной деформации, полученными И.И.\,Воровичем \cite{1} на 


основе гипотез Кирхгофа--Лява 


\begin{align} 


\varepsilon^0_{jj}&=w_{j\alpha^j}-B_{jj}w-G^k_{jj}w_k 


+\tfrac{1}{2}w^2_{\alpha^j},\quad j=1,2,\notag \\ 


\varepsilon^0_{12}&=\tfrac{1}{2}(w_{1\alpha^2}+w_{2\alpha^1})- 


B_{12}w-G^k_{12}w_k+\tfrac{1}{2}w_{\alpha^1}w_{\alpha^2}, \\ 


\varepsilon^1_{ij}&=-w_{\alpha^i\alpha^j}+G^k_{ij}w_{\alpha^k}, 


\quad i\le j, \ \; i,j=1,2,\notag 


\end{align} 


где $\varepsilon^0_{ij}$, $\varepsilon^1_{ij}$ --- компоненты 


тангенциальной и изгибной деформации 


срединной поверхности $S_0$; $w_i$, $w$ --- тангенциальные и 


нормальное 


перемещения точек $S_0$; $B_{ij}$ --- составляющие тензора кривизны 


$S_0$; $G^k_{ij}$ --- 


символы Кристоффеля второго рода; $\alpha^1$, $\alpha^2$ 


рассматриваются как декартовы 


координаты на плоскости, изменяющиеся в некоторой плоской области 


$\Omega$. 


 


Систему (1) будем решать относительно $w_i$, $w$, предполагая 


$\varepsilon^0_{ij}$, $\varepsilon^1_{ij}$ 


известными функциями. Для этого от (1) перейдем к системе 


\begin{gather} 


w_{1\alpha^1}-w_{2\alpha^2}+(B_{22}-B_{11})w+ 


(G^k_{22}-G^k_{11})w_k+\tfrac{1}{2}(w^2_{\alpha^1}- 


w^2_{\alpha^2})=\varepsilon^0_{11}-\varepsilon^0_{22}, \\ 


\begin{split} 


w_{1\alpha^2}+w_{2\alpha^1}-2B_{12}w-2G^k_{12}w_k+ 


w_{\alpha^1}w_{\alpha^2}&=2\varepsilon^0_{12},\\ 


-w_{\alpha^1\alpha^1}-w_{\alpha^2\alpha^2}+ 


(G^k_{11}+G^k_{22})w_{\alpha^k}&=\varepsilon^1_{11}+ 


\varepsilon^1_{22}. 


\end{split} 


\end{gather} 


Решение системы (2)-(3) начнем с решения уравнения (3). Введем 


обозначения: $a=-(G^1_{11}+G^1_{22})$, $b=G^2_{11}+G^2_{22}$, 


$\varepsilon_3=\varepsilon^1_{11}+\varepsilon^1_{22}$, 


$u=-w_{\alpha^1}$, $v=w_{\alpha^2}$. Тогда уравнение (3) будет 


эквивалентно системе 


\begin{gather*} 


u_{\alpha^1}-v_{\alpha^2}+au+bv=\varepsilon_3, \\ 


u_{\alpha^2}+v_{\alpha^1}=0, 


\end{gather*} 


которую с помощью комплексной функции 


$W(z)=u(\alpha^1,\alpha^2)+iv(\alpha^1,\alpha^2)$, 


$z=\alpha^1+i\alpha^2$ можно представить в виде 


\begin{equation} 


W_{\ov z}+AW+B\ov W=F, 


\end{equation} 


где $A=\tfrac{1}{4}(a-ib)$, $B=\tfrac{1}{4}(a+ib)$, 


$F=\varepsilon_3/2$, 


$W_{\ov z}=\tfrac{1}{2}(W_{\alpha^1}+iW_{\alpha^2})$. 


 


Уравнения вида (4)


изучены И.Н.\,Векуа \cite{3}. При нахождении решения уравнения (4) 


будем следовать \cite{3}. Предположим, что $\varepsilon_3\in 


L_2(\ov\Omega)$, функции $G^k_{ij}$, $B_{ij}$ ограничены 


в $\ov\Omega$; вне области $\ov\Omega$:


$\varepsilon_3=G^k_{ij}=B_{ij}\equiv 0$. При 


этих условиях решение (обобщенное) уравнения (4) можно получить в виде 


\begin{equation} 


W(z)=W_0(z)+(T_1\varepsilon_3)(z), 


\end{equation} 


где $W_0(z)$ --- обобщенная аналитическая функция \cite{3}, зависящая 


от произвольной голоморфной внутри $\Omega$ функции; оператор $T_1f$ 


дается формулой 


\begin{equation} 


T_1f=-\frac{1}{\pi}\intO \chi_1(z,\zeta;\Omega)f(\zeta) 


d\xi\,d\eta,\quad \zeta=\xi+i\eta, 


\end{equation} 


где $\chi_1(z,\zeta;\Omega)$ --- известная функция, называемая главной 


функцией \cite{3}. 


 


Зная $W(z)$, решение $w(z)$ уравнения (3) можно найти по формуле 


\begin{equation} 


w(\alpha^1,\alpha^2)=c_0-\re\int^z_{z_0}W(\zeta)d\zeta,\quad 


c_0=\const, 


\end{equation} 


$z_0$ --- произвольно фиксированная точка $\Omega$. 


 


Если $\Omega$ --- односвязная область, то правая часть (7) является 


однозначной функцией при фиксированных $c_0$, $z_0$. Если $\Omega$ --- 


многосвязная 


область, то правая часть (7) будет, вообще говоря, многозначной 


функцией. Пусть $\Gamma_j$ ($j=0,1,\dotsc,m$) --- замкнутые кривые, 


ограничивающие область 


$\Omega$, причем $\Gamma_1,\dotsc,\Gamma_m$ лежат внутри $\Gamma_0$. 


Тогда для однозначности 


правой части (7) необходимо и достаточно, чтобы выполнялись равенства 


$$ 


\re\int_{\Gamma_j}W(\zeta)d\zeta=0,\quad j=\ov{1,m}. 


$$ 


 


Перейдем к системе (2). При помощи комплексной функции $W_1=w_1+iw_2$ 


систему (2), как и выше, можно записать в виде 


\begin{equation} 


W_{1\,\ov z}+A_1W_1+B_1\ov W_1=F_1-K\varepsilon_3, 


\end{equation} 


где $F_1=\frac{1}{2}(\varepsilon_1+i\varepsilon_2)$, 


$\varepsilon_1=\varepsilon^0_{11}-\varepsilon^0_{22}$, 


$\varepsilon_2=2\varepsilon^0_{12}$; 


$K\varepsilon_3=\frac{1}{2}(K_1\varepsilon_3+iK_2\varepsilon_3)$ --- 


известный оператор, который определяется 


решением $w$ уравнения (3); коэффициенты $A_1$, $B_1$ зависят только от 


$G^k_{ij}$ и обращаются в нуль при $G^k_{ij}\equiv 0$. Предполагая, что 


$\varepsilon_1,\varepsilon_2\in L_2(\ov\Omega)$ и вне $\ov \Omega$:


$\varepsilon_i\equiv 0$, 


обобщенное решение уравнения (8) получаем в виде 


\begin{equation} 


W_1(z)=W_{1,0}(z)+T_2(\varepsilon_1-K_1\varepsilon_3)-T_3 


(\varepsilon_2-K_2\varepsilon_3), 


\end{equation} 


где через $W_{1,0}$ обозначена обобщенная аналитическая функция, 


зависящая 


от произвольной голоморфной внутри $\Omega$ функции; операторы 


$T_jf$ имеют ту же структуру (6), что и $T_1f$. 


 


\begin{lem} 


Операторы $T_jf$ $(j=\ov{1,3})$ суть вполне непрерывные линейные 


операторы в $L_2(\ov\Omega)$, 


отображающие это пространство в пространство 


$L_q(\ov\Omega)$, $q>1$, причем 


$$ 


\|T_jf\|_{L_q}\le c\|f\|_{L_2}. 


$$ 


\end{lem} 


 


Справедливость леммы 1 следует из свойств главных функций и теорем 


1.19, 1.27 (\cite{3}, сс.\,39,~46). 


 


Используя решения (5), (7), (9), их производные по $z$, $\ov z$, 


находим 


$w_{\alpha^j}$, $w_{\alpha^2\alpha^2}$, $w_{\alpha^1\alpha^2}$, 


$w_{2\alpha^2}$, с помощью которых затем из (1) получаем 


\begin{equation} 


\begin{split} 


\varepsilon^1_{22}&\equiv\varepsilon_4=H_1\varepsilon_3+ 


\tfrac{1}{2}\varepsilon_3+f_1,\quad 


\varepsilon^1_{12}\equiv\varepsilon_5=H_2\varepsilon_3+f_2, \\ 


\varepsilon^0_{22}&\equiv\varepsilon_6=H\varepsilon_3+ 


1_kH_{2+k}\varepsilon_k\big|_{k=\ov{1,3}}-\tfrac{1}{2} 


\varepsilon_1+f_3, 


\end{split} 


\end{equation} 


где $f_j\in L_q(\ov\Omega)$ ($q>1$) 


--- известные функции, зависящие от $W_0$, $W_{1,0}$; $H_jf$ --- 


известные операторы, определяемые, например, формулами 


\begin{gather} 


\begin{split} 


H_{2+k}\varepsilon_k&=\re[(-1)^k(S-I)P_{1+k}\varepsilon_k+ 


\tfrac{1}{2}i^{k-1}S\varepsilon_k]+(-1)^k\re 


(\ov g_2T_{1+k}\varepsilon_k),\quad k=1,2, \\ 


H\varepsilon_3&=\tfrac{1}{2}H_3[(\im T_1\varepsilon_3)^2- 


(\re T_1\varepsilon_3)^2]+H_4(\re T_1\varepsilon_3\cdot 


\im T_1\varepsilon_3)+\tfrac{1}{4}|T_1\varepsilon_3|^2, 


\end{split} 


\\ 


P_kf=A_1T_kf+B_1\ov{T_kf},\quad k=2,3;\quad 


g_k=G^1_{k,2}+iG^2_{k,2},\quad k=1,2, 


\end{gather} 


и т.\,д., $I$ --- тождественный оператор; оператор $Sf$ дается формулой 


$$


Sf=-\frac{1}{\pi}\intO f(\zeta)/(\zeta-z)^2d\xi\,d\eta,


$$


в которой интеграл следует понимать в смысле главного значения по Коши. 


Известно (\cite{3}, с.\,67), что $Sf$ 


есть ограниченный линейный оператор 


в $L_p(\ov\Omega)$, $p>1$, отображающий это пространство в себя. 


 


\begin{thm} 


$1)$ Операторы $H_jf$ $(j=\ov{1,4})$ суть линейные ограниченные 


операторы в $L_p(\ov\Omega)$, отображающие это пространство в себя, причем 


$\|H_jf\|_{L_p}\le c\|f\|_{L_p}$, $p>1$. 


 


$2)$ Оператор $H_5f$ --- линейный, $Hf$ 


--- нелинейный вполне непрерывный операторы в $L_2(\ov\Omega)$, 


отображающие это 


пространство в $L_q(\ov\Omega)$, $q>1$, причем 


$\|H_5f\|_{L_q}\le c\|f\|_{L_2}$, $\|Hf\|_{L_q}\le c\|f\|^2_{L_2}$, 


$q>1$. 


\end{thm} 


 


Теорема 1 доказывается с помощью леммы 1 и вышеприведенных свойств 


оператора $Sf$. 


 


Таким образом, получены выражения для компонент вектора перемещения 


точек срединной поверхности оболочки через $\varepsilon_1$, 


$\varepsilon_2$, $\varepsilon_3$ в виде (7), (9), 


которые будут использованы в дальнейшем. Шесть компонент деформации 


связаны между собой тремя соотношениями (10), которые представляют 


собой условия совместности деформации. При их выполнении решение 


системы (2)-(3) удовлетворяет системе (1). 


\medskip 


 


{\bf 2.} Пусть оболочка ортотропна, причем одна из осей ортотропии 


$\alpha^3$. Пусть в срединной поверхности $S_0$ существует ортогональная 


параметризация $x^1$, $x^2$, являющаяся параметризацией ортотропии. В 


главных осях ортотропии закон Гука запишется в виде (\cite{1}, с.\,31) 


$$ 


\varepsilon_{11}=\frac{\sigma_{11}}{E_1}- 


\frac{\nu_{12}\sigma_{22}}{E_2}-\frac{\nu_{13}\sigma_{33}}{E_3}, 


\quad \varepsilon_{12}=\frac{\sigma_{12}}{2G_{12}} \ \; 


(1\to 2\to 3\to 1). 


$$ 


Принимая гипотезы Кирхгофа--Лява, получаем 


$$ 


\sigma_{11}=\frac{E_1}{1-\nu_{12}\nu_{21}} 


(\varepsilon_{11}+\nu_{12}\varepsilon_{22}),\quad 


\sigma_{12}=2G_{12}\varepsilon_{12} \ \; (1\rightleftarrows 2). 


$$ 


Тогда для объемной плотности потенциальной энергии деформации оболочки 


$\Pi$ имеем формулу 


\begin{equation} 


2\Pi=B^{\lambda\mu q s}(x,\alpha^3)\gamma_{\lambda\mu} 


(x,\alpha^3)\gamma_{q s}(x,\alpha^3),\quad 


\lambda\le\mu,\ \; q\le s, 


\end{equation} 


где $B^{1111}=\frac{E_1}{1-\nu_{12}\nu_{21}}$, 


$B^{1122}=\frac{1}{2} 


\frac{(E_1\nu_{12}+E_2\nu_{21})}{1-\nu_{12}\nu_{21}}\; 


(1\rightleftarrows 2)$, $B^{1212}=G_{12}$; 


$\gamma_{jj}=\varepsilon_{jj}$, 


$\gamma_{12}=2\varepsilon_{12}$. 


 


{\it Условие} А: {\it предполагается, что $E_i$, $G_{12}$, $\nu_{ij}$ 


--- ограниченные функции переменных $\alpha^1$, $\alpha^2$, 


$\alpha^3$}; квадратичная форма $2\Pi$ 


равномерно-положительно определена во всем объеме $V$, занятом 


оболочкой, т.\,е. имеет место неравенство (\cite{1}, с.\,32) 


\begin{equation} 


2\Pi\ge c(\gamma^2_{11}+\gamma^2_{12}+\gamma^2_{22}),\quad c>0. 


\end{equation} 


 


Для потенциальной энергии $U$ деформации, накопленной во всем объеме 


оболочки, с учетом тонкостенности оболочки, формулы (13) и 


$\gamma_{ij}=\gamma^0_{ij}+\alpha^3\gamma^1_{ij}$ \cite{1}, 


будем иметь 


\begin{equation} 


U=\intO QD d\alpha^1\,d\alpha^2, 


\end{equation} 


где 


\begin{gather} 


2Q=D^{\lambda\mu q s}_p\gamma^0_{\lambda\mu}\gamma^0_{q s}+ 


D^{\lambda\mu q s}_* 


(\gamma^1_{\lambda\mu}\gamma^0_{q s}+ 


\gamma^0_{\lambda\mu}\gamma^1_{q s})+ 


D_u^{\lambda\mu q s}\gamma^1_{\lambda\mu} 


\gamma^1_{q s}, \\ 


%


D^{\lambda\mu q s}_p(x)=\int^h_{-h}B^{\lambda\mu q s} 


(x,\alpha^3)d\alpha^3,\ \;


D^{\lambda\mu q s}_*(x)=\int^h_{-h}B^{\lambda\mu q s} 


(x,\alpha^3)\alpha^3d\alpha^3,\notag \\


%


D^{\lambda\mu q s}_u(x)=\int^h_{-h} 


B^{\lambda\mu q s}(x,\alpha^3)(\alpha^3)^2d\alpha^3,\ \;


\lambda\le \mu,\ \;q\le s; \notag


\end{gather} 


$D(\alpha^1,\alpha^2)$ --- элемент площади срединной поверхности. 


Предположим, что в области $\ov\Omega$ 


\begin{equation} 


D(\alpha^1,\alpha^2)\ge c^2>0. 


\end{equation} 


 


В произвольных $s$-координатах $\alpha^1$, $\alpha^2$, $\alpha^3$ 


потенциальная энергия дается 


формулой (15), где $Q$ определяется по формуле (16) с коэффициентами 


$D^{\lambda\mu q s}_{p,*,u}(\alpha)$, 


которые известным образом (\cite{1}, с.\,34) выражаются через 


$D^{\lambda\mu q s}_{p,*,u}(x)$. Легко 


видеть, что в силу (14) квадратичная форма $2Q$ является положительно 


определенной формой. Отсюда следует положительная определенность 


тензоров $D^{\lambda\mu q s}_{p,u}$. 


 


В (16) перейдем к переменным $\varepsilon_i$ ($i=\ov{1,6}$). Из 


рассуждений п.\,1 следует 


\begin{equation} 


\gamma^0_{11}=\varepsilon_1+\varepsilon_6,\quad 


\gamma^0_{12}=\varepsilon_2,\quad \gamma^0_{22}=\varepsilon_6, 


\quad \gamma^1_{11}=\varepsilon_3-\varepsilon_4,\quad 


\gamma^1_{12}=2\varepsilon_5,\quad \gamma^1_{22}=\varepsilon_4. 


\end{equation} 


Соотношения (18) представляют собой невырожденные преобразования 


переменных. Поэтому $2Q$ относительно новых переменных $\varepsilon_i$ 


также является 


положительно определенной квадратичной формой, т.\,е. имеет место 


неравенство 


\begin{equation} 


2Q\ge c(\varepsilon^2_1+\varepsilon^2_2+\varepsilon^2_3+ 


\varepsilon^2_4+\varepsilon^2_5+\varepsilon^2_6),\quad c>0. 


\end{equation} 


Используя формулы (18), условия совместности деформации (10), 


$\gamma^0_{ij}$, $\gamma^1_{ij}$ представим в виде 


\begin{equation} 


\gamma^0_{jj}=t^0_{jj}+H\varepsilon_3+f_3,\quad 


\gamma^1_{jj}=t^1_{jj}+(-1)^jf_1,\quad j=1,2;\quad 


\gamma^0_{12}=t^0_{12},\quad \gamma^1_{12}=t^1_{12}+2f_2, 


\end{equation} 


где через $t^0_{ij}$, $t^1_{ij}$ обозначены линейные части 


$\gamma^0_{ij}$, $\gamma^1_{ij}$: 


\begin{equation} 


\begin{split} 


t^0_{11}&=\tfrac{1}{2}\varepsilon_1+1_kH_{2+k} 


\varepsilon_k\big|_{k=\ov{1,3}},\quad t^0_{12}= 


\varepsilon_2,\quad t^0_{22}=1_kH_{2+k} 


\varepsilon_k\big|_{k=\ov{1,3}}-\tfrac{1}{2}\varepsilon_1, \\ 


t^1_{11}&=\tfrac{1}{2}\varepsilon_3-H_1\varepsilon_3,\quad 


t^1_{12}=2H_2\varepsilon_3,\quad 


t^1_{22}=H_1\varepsilon_3+\tfrac{1}{2}\varepsilon_3 


\end{split} 


\end{equation} 


Теперь в соответствии с формулами (20) для вариации $2Q$ 


получаем формулу 


\begin{equation} 


\delta Q=D^{\lambda\mu q s}_p t^0_{\lambda\mu}\delta t^0_{qs} 


+D^{\lambda\mu qs}_*(t^0_{\lambda\mu} 


\delta t^1_{qs}+t^1_{\lambda\mu}\delta t^0_{qs})+ 


D^{\lambda\mu qs}_u t^1_{\lambda\mu}\delta t^1_{qs} 


+Q_p(\varepsilon;\delta\varepsilon)+Q_*(\varepsilon;\delta\varepsilon) 


+\Phi(\delta\varepsilon), 


\end{equation} 


где 


\begin{gather} 


Q_p(\varepsilon;\delta\varepsilon)=[d(H\varepsilon_3+f_3)+ 


d^{\lambda\mu}_p t^0_{\lambda\mu}(\varepsilon)]H_\delta 


(\varepsilon_3,\delta\varepsilon_3)+H\varepsilon_3


d^{\lambda\mu}_p t^0_{\lambda\mu}(\delta\varepsilon), \\ 


%


Q_*(\varepsilon;\delta\varepsilon)=d^{\lambda\mu}_* 


[t^1_{\lambda\mu}(\varepsilon)+\beta_{\lambda\mu}]H_\delta 


(\varepsilon_3,\delta\varepsilon_3)+H\varepsilon_3 


d^{\lambda\mu}_* t^1_{\lambda\mu}(\delta\varepsilon), \\ 


%


\delta H\varepsilon_3\equiv H_\delta(\varepsilon_3, 


\delta\varepsilon_3),\quad 


\varepsilon=(\varepsilon_1,\varepsilon_2,\varepsilon_3),\quad 


\delta\varepsilon=(\delta\varepsilon_1,\delta\varepsilon_2, 


\delta\varepsilon_3); \notag


\end{gather} 


$d^{\lambda\mu}_{p,*}$, $d$, $\beta_{ij}$ --- известные функции, 


зависящие от $D^{\lambda\mu qs}_{p,*}$ и $f_k$ соответственно; 


через $\Phi(\delta\varepsilon)$ обозначено известное выражение, 


содержащее только вариацию $\varepsilon$. 


 


Вычислим работу внешних приложенных к оболочке сил на возможных перемещениях


в условиях гипотез Кирхгофа--Лява. Пусть по граням $S^\pm_0$ оболочки 


приложены 


усилия $\bold F^\pm(\alpha^1,\alpha^2)$, действуют массовые силы 


$\bold F(\alpha^1,\alpha^2,\alpha^3)$ и на 


границе $S^0_0$ действуют 


поверхностные силы $\bold F(s,\alpha^3)$. Предположим, 


что внешние силы таковы, 


что движение оболочки как твердого тела отсутствует. В связи с этим в 


дальнейшем будем считать, что произвольные функции $W_0$, $W_{1,0}$ 


равны нулю. 


 


Для элементарной работы $\delta A_1$ массовых сил и сил, приложенных к 


$S^\pm_0$, имеем 


$$ 


\delta A_1=\intO R^j\delta w_j\big|_{j=\ov{1,3}} 


D d\alpha^1 d\alpha^2-\intO \cal L^i\delta w_{\alpha^i} 


\big|_{i=1,2}D d\alpha^1d\alpha^2. 


$$ 


Элементарная работа $\delta A_2$ усилий, приложенных к $S^0_0$, дается 


формулой 


$$ 


\delta A_2=\int_\Gamma(Q^j\delta w_j\big|_{j=\ov{1,3}}- 


M^i\delta w_{\alpha^i}\big|_{i=1,2})ds. 


$$ 


Здесь $R^j$, $\cal L^i$, $Q^j$, $M^i$ --- известные функции. Отметим, 


если 


{\it условие}~B\; $\bold F\in L_2(\ov\Omega)\times L_1[-h,h]$, 


$\bold F^\pm\in L_2(\ov\Omega)$, то 


$R^j,\cal L^i\in L_2(\ov\Omega)$, $Q^j,M^i\in L_2(\Gamma)$. Принимая во 


внимание соотношения 


(5), (7), (9), (10), выражения для $\delta A_j$ можно преобразовать к 


виду 


\begin{align} 


\delta A_1&=\intO D[R^1\re T_4(\varepsilon_3,\delta\varepsilon_3) 


+R^2\im T_4(\varepsilon_3,\delta\varepsilon_3)] 


d\alpha^1d\alpha^2+I_1(\delta\varepsilon), \\ 


\delta A_2&=\intO[\wt Q^2K^0_2(\varepsilon_3,\delta\varepsilon_3)- 


\wt Q^1K^0_1(\varepsilon_3,\delta\varepsilon_3)] 


d\alpha^1d\alpha^2+I_2(\delta\varepsilon), 


\end{align} 


где 


\begin{align} 


T_4(\varepsilon_3,\delta\varepsilon_3)&= 


T_3 K^0_2(\varepsilon_3,\delta\varepsilon_3)- 


T_2K^0_1(\varepsilon_3,\delta_3),\notag \\ 


K^0_1(\varepsilon_3,\delta\varepsilon_3)&= 


\re T_1\varepsilon_3\,\re T_1\delta\varepsilon_3- 


\im T_1\varepsilon_3\,\im T_1\delta\varepsilon_3, \\ 


K^0_2(\varepsilon_3,\delta\varepsilon_3)&= 


(-\tfrac{1}{2})[\re T_1\varepsilon_3\,\im T_1\delta\varepsilon_3+ 


\im T_1\varepsilon_3\,\re T_1\delta\varepsilon_3];\notag 


\end{align} 


через $I_j(\delta\varepsilon)$ обозначены линейные ограниченные 


относительно $\delta\varepsilon$ функционалы в 


$L_2(\ov\Omega)$; $\wt Q^j$ ($\in L_2(\ov\Omega)$) --- известные 


функции. 





Введем линейное пространство $\wt L_2(\ov\Omega)$, 


любой паре элементов 


$\varepsilon^j=(\varepsilon^j_1,\varepsilon^j_2,\varepsilon^j_3)$ 


которого поставлено в соответствие число 


\begin{multline} 


(\varepsilon^1,\varepsilon^2)_{\wt L_2}=\intO 


\{D^{\lambda\mu qs}_p t^0_{\lambda\mu}(\varepsilon^1) 


t^0_{qs}(\varepsilon^2)+ \\ 


%


+ D^{\lambda\mu qs}_*[t^0_{\lambda\mu}(\varepsilon^1) 


t^1_{qs}(\varepsilon^2)+t^1_{\lambda\mu}(\varepsilon^1) 


t^0_{qs}(\varepsilon^2)]+D^{\lambda\mu qs}_u 


t^1_{\lambda\mu}(\varepsilon^1)t^1_{qs}(\varepsilon^2)\} 


D d\alpha^1d\alpha^2, 


\end{multline} 


удовлетворяющее всем условиям скалярного произведения. Проверим лишь, 


что из $(\varepsilon,\varepsilon)_{\wt L_2}=0$ следует $\varepsilon=0$ 


(справедливость остальных условий очевидна). 


Действительно, в силу положительной определенности квадратичной формы 


$2Q$ с учетом (17) из $(\varepsilon,\varepsilon)_{\wt L_2}=0$ следует 


$t^0_{\lambda\mu}(\varepsilon)=t^1_{\lambda\mu}(\varepsilon)=0$. Тогда, 


принимая во внимание формулы (21), сразу получаем $\varepsilon=0$. 


 


Скалярному произведению (28) соответствует норма 


\begin{multline} 


\|\varepsilon\|^2_{\wt L_2}=\intO\{D^{\lambda\mu qs}_p 


t^0_{\lambda\mu}(\varepsilon)t^0_{qs}(\varepsilon)+ \\ 


+D^{\lambda\mu qs}_*[t^0_{\lambda\mu}(\varepsilon) 


t^1_{qs}(\varepsilon)+t^1_{\lambda\mu}(\varepsilon) 


t^0_{qs}(\varepsilon)]+D^{\lambda\mu qs}_u t^1_{\lambda\mu} 


(\varepsilon)t^1_{qs}(\varepsilon)\}D d\alpha^1d\alpha^2. 


\end{multline} 


 


\begin{lem} 


Пусть выполнено условие {\em A}. Тогда имеет место 


$$ 


m_0\|\varepsilon\|^2_{L_2}\le\|\varepsilon\|^2_{\wt L_2}\le 


M_0\|\varepsilon\|^2_{L_2},\quad m_0,M_0>0, 


$$ 


т.\,е. нормы $\|\cdot\|_{L_2}$, $\|\cdot\|_{\wt L_2}$ эквивалентны. 


\end{lem} 


 


Справедливость леммы 2 следует из неравенства (19) и условия A. 


 


Из леммы 2 следует, что пространство 


$\wt L_2(\ov\Omega)$ с нормой (29) является 


гильбертовым пространством. 





Введем понятие обобщенного решения задачи в деформациях.


 


\begin{defnnonum} 


Обобщенным решением задачи в деформациях назовем вектор-функцию 


$\varepsilon=(\varepsilon_1,\varepsilon_2,\varepsilon_3)\in 


\wt L_2(\ov\Omega)$, 


удовлетворяющую интегральному соотношению 


\begin{multline} 


(\varepsilon,\varphi)_{\wt L_2}=\intO\{D[R^1\re T_4 


(\varepsilon_3,\varphi_3)+R^2\im T_4(\varepsilon_3,\varphi_3)] 


+\wt Q^2K^0_2(\varepsilon_3,\varphi_3)- 


\wt Q^1K^0_1(\varepsilon_3,\varphi_3)\}d\alpha^1d\alpha^2- \\ 


-\intO D[Q_p(\varepsilon,\varphi)+ 


Q_*(\varepsilon,\varphi)]d\alpha^1d\alpha^2- I(\varphi) 


\end{multline} 


для любой вектор-функции 


$\varphi=(\varphi_1,\varphi_2,\varphi_3)\in\wt L_2(\ov\Omega)$. 


Здесь $I(\varphi)$ --- линейный 


функционал в 


$\wt L_2(\ov\Omega)$; $Q_{p,*}(\varepsilon,\varphi)$, 


$T_4(\varepsilon_3,\varphi_3)$, $K^0_j(\varepsilon_3,\varphi_3)$ 


даются формулами (23), (24), (27) соответственно. 


\end{defnnonum} 


 


Отметим, что при введении обобщенного решения задачи в виде (30) мы 


следовали \cite{1}. Если принять во внимание соотношения (15), (22), 


(20), (25), (26), то легко видеть, что (30) выражает принцип Лагранжа 


для системы оболочка--внешние силы. При этом 


$\varphi=(\varphi_1,\varphi_2,\varphi_3)$ 


суть возможные деформации для этой системы. 


 


\begin{lem} 


Пусть выполнены условия {\em A} и {\em B}. Тогда при 


$\varepsilon,\varphi\in\wt L_2(\ov\Omega)$ интегралы $I(\varphi)$ 


в $(30)$ имеют смысл и представляют собой линейные ограниченные 


функционалы в $\wt L_2(\ov\Omega)$ 


относительно $\varphi$ при фиксированной вектор-функции $\varepsilon$. 


\end{lem} 


 


Прежде всего отметим, что в силу леммы 2\; 


$\varepsilon,\varphi\in L_2(\ov\Omega)$. Далее доказательство 


леммы 3 проводится с помощью формул (27), леммы 1, теоремы 1. 


 


По теореме Рисса существует элемент 


$G\in\wt L_2(\ov\Omega)$ такой, что правая часть (30) 


представляется в виде скалярного произведения 


$(G,\varphi)_{\wt L_2}$. Тогда интегральное 


соотношение (30) запишется в виде 


$(\varepsilon,\varphi)_{\wt L_2}=(G,\varphi)_{\wt L_2}$. Отсюда в силу 


произвольности $\varphi$ получаем 


\begin{equation} 


\varepsilon-G(\varepsilon)=0. 


\end{equation} 


Таким образом, нахождение обобщенного решения задачи свелось к 


нахождению решения операторного уравнения (31). 





Займемся исследованием разрешимости уравнения (31).


 


\begin{thm} 


Оператор $G(\varepsilon)$ действует из 


$\wt L_2(\ov\Omega)$ в $\wt L_2(\ov\Omega)$ усиленно непрерывно. 


\end{thm} 


 


Действительно, пусть $\varepsilon^n\to\varepsilon^0$ слабо в 


$\wt L_2(\ov\Omega)$. В силу 


леммы 2 можем считать, что 


$\varepsilon^n\to\varepsilon^0$ слабо в 


$L_2(\ov\Omega)$. Так как $T_jf$ ($j=\ov{1,3}$), $H_5f$ --- вполне 


непрерывные линейные 


операторы в $L_2(\ov\Omega)$, 


то они усиленно непрерывны (\cite{1}, с.\,65). Далее, в 


силу усиленной непрерывности $T_1f$ и ограниченности операторов $H_3f$, 


$H_4f$ в $L_2(\ov\Omega)$ 


из формулы (11) следует усиленная непрерывность операторов $Hf$, 


$H_\delta(\varepsilon_3,\varphi_3)$ (относительно $\varepsilon_3$). 


Тогда, используя соотношения (23), (24), 


(21), (27), после громоздких выкладок получим оценку 


$$ 


|(G(\varepsilon^n)-G(\varepsilon^0),\varphi)_{\wt L_2}| 


\le \delta_n\|\varphi\|_{\wt L_2}, 


$$ 


где $\delta_n\to 0$ при $n\to\infty$. Отсюда следует, что 


$\|G(\varepsilon^n)-G(\varepsilon^0)\|_{\wt L_2}\le\delta_n$, т.\,е. 


оператор $G(\varepsilon)$ усиленно 


непрерывен в $\wt L_2(\ov\Omega)$, 


следовательно, вполне непрерывен (\cite{1}, с.\,65). 


 


\begin{lem} 


Пусть $\Gamma\in C^1_\alpha$ $(0<\alpha<1)$ и выполнено условие 


\begin{equation} 


q\equiv q_0\|\ov{\Phi_\Gamma(z)}+2(\re T^*_2\gamma- 


i\re T^*_3\gamma)\|_C<1, 


\end{equation} 


где 


$\gamma=A_1+\ov B_1-\ov g_2-\Phi_\Gamma A_1-\ov{\Phi_\Gamma B_1};$ 


$\Phi_\Gamma(z)=\int\limits_\Gamma 1/(z-\zeta)d\ov\zeta 


\frac{1}{2\pi i};$ $q_0$ 


$(\le 1)$ --- известная положительная постоянная$;$ $T^*_jf$ 


--- оператор, сопряженный оператору $T_jf$. Тогда уравнение 


\begin{equation} 


(H\varepsilon_3)(z)=0,\quad x\in\ov\Omega 


\end{equation} 


имеет только тривиальное решение в $L_2(\ov\Omega)$. 


\end{lem} 


 


Покажем это. Пусть 


$\varepsilon_3\in L_2(\ov\Omega)$ есть нетривиальное решение уравнения 


(33). Обе части (33) интегрируем по области $\Omega$, при этом 


используем формулы (11), (12). Применяя формулу Грина (\cite{3}, 


с.\,28) к интегралам от $Sf$ по 


области $\Omega$ и меняя порядок интегрирования, после несложных 


преобразований приходим к неравенству вида 


$$ 


\intO|T_1\varepsilon_3|^2d\alpha^1d\alpha^2\le 


q\intO|T_1\varepsilon_3|^2d\alpha^1d\alpha^2. 


$$ 


Отсюда в силу (32) следует, что $T_1\varepsilon_3\equiv 0$ в 


$\ov\Omega$. Так 


как $g(z)=T_1\varepsilon_3$ удовлетворяет уравнению (4), имеем 


$\varepsilon_3\equiv 0$ в $\ov\Omega$. Лемма доказана. 


 


Отметим, если $\Omega$ --- круг, то 


$\Phi_\Gamma(z)\equiv 0$, $z\in\Omega$. Кроме этого, если срединная 


поверхность $S_0$ является поверхностью нулевой гауссовой кривизны, то, 


как известно (\cite{1}, с.\,18), на $S_0$ существует единая евклидова 


параметризация. Тогда $G^k_{ij}\equiv 0$, следовательно, 


$A_1=B_1=g_2\equiv 0$. 


Отсюда $q=0$, т.\,е. условие (32) в этом частном случае выполняется. 


 


В дальнейшем, как и в \cite{1}, потребуется некоторое разбиение сферы 


$S_{\wt L_2}(R,0)$ гильбертова пространства 


$\wt L_2(\ov\Omega)$ радиуса $R$ с центром в 


начале координат. 


Для его построения возьмем сферу $\|\varepsilon_3\|_{L_2}=1$ 


(обозначим ее через $S_{L_2}(1,0)$) 


пространства $L_2(\ov\Omega)$. 


Через $S'_{L_2}(1,0)$ обозначим подмножество $S_{L_2}(1,0)$, элементы 


которого удовлетворяют неравенству 


\begin{equation} 


\negthickspace\negthickspace\negthickspace


I(\varepsilon_3)\equiv\tfrac{c_1}{\delta} 


\|H_5\varepsilon_3\|^2_{L_2}\!+\!c_2\|\beta_1\re 


T_4(\varepsilon_3,\varepsilon_3)\!+\!\beta_2\im


T_4(\varepsilon_3,\varepsilon_3)\!+ \!


\beta_3 K^0_2(\varepsilon_3,\varepsilon_3)\!- \!


\beta_4K^0_1(\varepsilon_3,\varepsilon_3)\|_{L_2}> 


\tfrac{1}{2}, 


\end{equation} 


где $c_j$, $\delta>0$ --- некоторые фиксированные постоянные; $\beta_j$ 


--- фиксированные 


функции класса $L_2(\ov\Omega)$, 


не зависящие от $\varepsilon_3$. Множество $S'_{L_2}(r_3,0)$ в силу 


однородности 


функционала $I(\varepsilon_3)$ есть центральная проекция 


$S'_{L_2}(1,0)$ с единичной сферы 


$S_{L_2}(1,0)$ на сферу $S_{L_2}(r_3,0)$. Пусть $S''_{L_2}(1,0)$ есть 


подмножество $S_{L_2}(1,0)$, элементы которого удовлетворяют условию 


\begin{equation} 


I(\varepsilon_3)\le\frac{1}{2}. 


\end{equation} 


$S''_{L_2}(r_3,0)$ есть центральная проекция $S''_{L_2}(1,0)$ с 


$S_{L_2}(1,0)$ на $S_{L_2}(r_3,0)$. Очевидно, 


$S'_{L_2}(r_3,0)\cup S''_{L_2}(r_3,0)=S_{L_2}(r_3,0)$. Если одно из 


подмножеств $S'_{L_2}(r_3,0)$, $S''_{L_2}(r_3,0)$ окажется пустым, это 


лишь упростит наши рассуждения. Через $\ov S\,{}'_{L_2}(1,0)$ обозначим 


слабое замыкание $S'_{L_2}(1,0)$. 


 


\begin{lem} 


Множество $\ov S\,{}'_{L_2}(1,0)$ не содержит нуля. 


\end{lem} 


 


Действительно, т.\,к. $H_50=T_4(0,0)=K^0_j(0,0)=0$, множество 


$S'_{L_2}(1,0)$ не содержит нуля. Пусть 


$\varepsilon^n_3\in S'_{L_2}(1,0)$ и $\varepsilon^n_3\to 0$ 


(слабо) в $L_2(\ov\Omega)$, $n\to\infty$. 


Из усиленной непрерывности операторов $H_5\varepsilon_3$, 


$T_4(\varepsilon_3,\varepsilon_3)$, 


$K^0_j(\varepsilon_3,\varepsilon_3)$ 


в $L_2(\ov\Omega)$ следует 


$$ 


\|H_5\varepsilon^n_3\|_{L_2},\, 


\|T_4(\varepsilon^n_3,\varepsilon^n_3)\|_{L_2}, \, 


\|K^0_j(\varepsilon^n_3,\varepsilon^n_3)\|_{L_2} 


\longrightarrow 0\quad (n\to \infty) 


$$ 


и (34) становится невозможным. Полученное противоречие доказывает 


лемму 5. 


 


\begin{thm} 


Пусть выполнено условие $(32)$. Тогда на $S'_{L_2}(r_3,0)$ справедливо 


неравенство 


$$ 


\|H\varepsilon_3\|^2_{L_2}\ge c\,r^4_3,\quad c>0. 


$$ 


\end{thm} 


 


Вначале установим, что на $S'_{L_2}(1,0)$ имеет место 


\begin{equation} 


\|H\varepsilon_3\|^2_{L_2}\ge c>0. 


\end{equation} 


Пусть существует последовательность 


$\varepsilon^n_3\in\ov S\,{}'_{L_2}(1,0)$ такая, что 


$\|H\varepsilon^n_3\|_{L_2}\to 0$ ($n\to\infty$). Предположим, 


что $\varepsilon^n_3\to\varepsilon^0_3$ (слабо), 


$\varepsilon^0_3\in\ov S\,{}'_{L_2}(1,0)$. Так как оператор 


$H\varepsilon_3$ усиленно непрерывен в $L_2(\ov\Omega)$, то 


$H\varepsilon^n_3\Longrightarrow H\varepsilon^0_3$ 


(сильно) в $L_2(\ov\Omega)$. 


Следовательно, $H\varepsilon^0_3=0$. Тогда в силу леммы 4\; 


$\varepsilon^0_3\equiv 0$, 


что противоречит лемме 5. Итак, на $S'_{L_2}(1,0)$ имеет место (36). 


Отсюда в силу однородности оператора $H\varepsilon_3$ сразу получаем 


утверждение теоремы 3. 


 


Введем функционал 


$N(\varepsilon,t)=(\varepsilon-t G(\varepsilon), 


\varepsilon)_{\wt L_2}$, 


определенный на $\wt L_2(\ov\Omega)\times[0,1]$. 


Для исследования этого функционала сначала 


$Q_p(\varepsilon,\varepsilon)$ преобразуем к виду 


\begin{equation} 


Q_p(\varepsilon,\varepsilon)=2Q^0_p(\varepsilon,\varepsilon)- 


H\varepsilon_3(d^{12}_p t^0_{12}-3d^{11}_p t^0_{11}- 


3d^{22}_p t^0_{22}-2d f_3)- 


2D^{1212}_p(t^0_{12})^2, 


\end{equation} 


где 


$Q^0_p(\varepsilon,\varepsilon)=D^{\lambda\mu qs}_p 


a_{\lambda\mu}a_{qs}$, $a_{11}=a_{22}=H\varepsilon_3$, 


$a_{12}=t^0_{12}(\varepsilon)$. В 


силу положительной определенности $D^{\lambda\mu qs}_p$ имеет место 


неравенство 


\begin{equation} 


Q^0_p(\varepsilon,\varepsilon)\ge c 


[2(H\varepsilon_3)^2+(t^0_{12})^2],\quad c>0. 


\end{equation} 


Предположим, что выполняется условие 


\begin{equation} 


\intO D[Q^0_p(\varepsilon,\varepsilon)+ 


Q_*(\varepsilon,\varepsilon)]d\alpha^1d\alpha^2\ge 0, 


\end{equation} 


где $Q_*(\varepsilon,\varepsilon)$ дается формулой (24). 


 


\begin{thm} 


Пусть выполнены условия {\em A, B, (32), (39)}. Тогда на сферах 


$S_{\wt L_2}(R,0)$ 


достаточно большого радиуса $R$ справедливо неравенство 


$$ 


N(\varepsilon,t)\ge c\,R^2,\quad 


c>0,\ \; 0\le t\le 1. 


$$ 


\end{thm} 


 


Покажем это. Через $S'_{\wt L_2}(R,0)$, $S''_{\wt L_2}(R,0)$ обозначим 


множества вектор-функций 


$\varepsilon=(\varepsilon_1,\varepsilon_2,\varepsilon_3)\in 


S_{\wt L_2}(R,0)$, 


у которых $\varepsilon_3$ принадлежит $S'_{L_2}(r_3,0)$, 


$S''_{L_2}(r_3,0)$ соответственно. 


Очевидно, 


$S_{\wt L_2}(R,0)=S'_{\wt L_2}(R,0)\cup S''_{\wt L_2}(R,0)$. 


Оценим вначале $N(\varepsilon,t)$ на $S'_{\wt L_2}(R,0)$. 


С учетом (37), (39) будем иметь 


\begin{multline} 


N(\varepsilon,t)\ge\|\varepsilon\|^2_{\wt L_2}+ 


t\intO DQ^0_p(\varepsilon,\varepsilon)d\alpha^1d\alpha^2- 


t\intO D[|H\varepsilon_3(d^{12}_p t^0_{12}- 


3d^{11}_p t^0_{11}-3d^{22}_p t^0_{22})|+ \\ 


%


+D^{1212}_p(t^0_{12})^2]d\alpha^1d\alpha^2- 


t\biggl|\intO\{D[2df_3H\varepsilon_3- 


R^1\re T_4(\varepsilon_3,\varepsilon_3)-R^2\im T_4 


(\varepsilon_3,\varepsilon_3)]\biggr.- \\ 


%


-\biggl.\wt Q^2K^0_2(\varepsilon_3,\varepsilon_3)+ 


\wt Q^1K^0_1(\varepsilon_3,\varepsilon_3)\}\biggr| 


d\alpha^1d\alpha^2-t| I(\varepsilon)|. 


\end{multline} 


Так как $I(\varepsilon)$ есть линейный ограниченный функционал в 


$\wt L_2(\ov\Omega)$, то существует 


$\varepsilon^0\in\wt L_2(\ov\Omega)$ 


такое, что $I(\varepsilon)=(\varepsilon^0,\varepsilon)_{\wt L_2}$. 


Далее, используя (17), (38), теоремы 1, 3, из (40) будем иметь 


\begin{multline} 


N(\varepsilon,t)\ge \|\varepsilon\|^2_{\wt L_2}+tc[c 


\|\varepsilon_3\|^4_{L_2}-\|\varepsilon_3\|^3_{L_2}- \\


-(1+\|\varepsilon_1\|_{L_2}+\|\varepsilon_2\|_{L_2}) 


\|\varepsilon_3\|^2_{L_2}-\|\varepsilon^0\|_{\wt L_2} 


\|\varepsilon_3\|_{L_2}-\|\varepsilon_2\|^2_{L_2}], \ \; c>0. 


\end{multline} 


Выражение в квадратных скобках представляет собой полином четвертой 


степени относительно $\|\varepsilon_3\|_{L_2}$ с положительным старшим 


коэффициентом. Поэтому 


при произвольно фиксированных 


$\|\varepsilon_j\|_{L_2}=r_j$ ($j=1,2$) радиус сферы 


$\|\varepsilon_3\|_{L_2}=r_3$ можем взять 


настолько большим, что выражение в квадратных скобках в (41) будет 


больше нуля. Тогда в силу леммы 2 при


$R^2\in[m_0r^2,M_0r^2]$, $r^2=r^2_1+r^2_2+r^2_3$ на 


$S'_{\wt L_2}(R,0)$ имеет место 


неравенство 


\begin{equation} 


N(\varepsilon,t)\ge R^2,\quad 0\le t\le 1. 


\end{equation} 


Теперь оценим $N(\varepsilon,t)$ на $S''_{\wt L_2}(R,0)$. 


Из (40) имеем 


\begin{multline} 


N(\varepsilon,t)\ge\|\varepsilon\|^2_{\wt L_2}+t\biggl[\, 


\intO DQ^0_p(\varepsilon,\varepsilon)d\alpha^1d\alpha^2- 


\delta\|H\varepsilon_3\|^2_{L_2}\biggr]- \\ 


%


-t\biggl\{\frac{1}{2\delta}\|D(d^{12}_p t^0_{12}- 


3d^{11}_p t^0_{11}-3d^{22}_p t^0_{22})\|^2_{L_2}+ 


2\|DD^{1212}_p\|_c\,\|t^0_{12}\|^2_{L_2}+ 


\frac{1}{\delta}\|df_3D\|^2_{L_2}\biggr\}-\\


%


-\sqrt{\mes \Omega}\|D[R^1\re T_4(\varepsilon_3,\varepsilon_3)\!+\!


R^2\im T_4(\varepsilon_3,\varepsilon_3)]\!+ \!


\wt Q^2K^0_2(\varepsilon_3,\varepsilon_3)\!- \!


\wt Q^1K^0_1(\varepsilon_3,\varepsilon_3)\|_{L_2}\!- \!


\|\varepsilon^0\|_{\wt L_2}\,\|\varepsilon\|_{\wt L_2}. 


\end{multline} 


В силу того, что $Q^0_p(\varepsilon,\varepsilon)$ является положительно 


определенной 


формой переменных $H\varepsilon_3$, $t^0_{12}$, постоянную $\delta>0$ 


можно выбрать так, чтобы 


\begin{equation} 


\intO DQ^0_p(\varepsilon,\varepsilon)d\alpha^1d\alpha^2- 


\delta\|H\varepsilon_3\|^2_{L_2}>0. 


\end{equation} 


Теперь в силу (44) из (43) получаем 


\begin{multline} 


N(\varepsilon,t)\ge\|\varepsilon\|^2_{\wt L_2}-\frac{3}{2\delta} 


\|dD\|_c\,\|H_5\varepsilon_3\|^2_{L_2}-\frac{c}{2\delta} 


(\|\varepsilon_1\|^2_{L_2}+\|\varepsilon_2\|^2_{L_2})- \\


- 2\|DD^{1212}_p\|_c\,\|\varepsilon_2\|^2_{L_2}-\frac{1}{\delta} 


\|df_3D\|^2_{L_2} - 


\sqrt{\mes \Omega}\|D[R^1\re T_4(\varepsilon_3,\varepsilon_3)+ 


R^2\im T_4(\varepsilon_3,\varepsilon_3)]+ \\


+\wt Q^2K^0_2(\varepsilon_3,\varepsilon_3)- 


\wt Q^1K^0_1(\varepsilon_3,\varepsilon_3)\|_{L_2}- 


\|\varepsilon^0\|_{\wt L_2}\,\|\varepsilon\|_{\wt L_2}.


\end{multline} 


Положим $C_1=\frac{3}{2m_0}\|dD\|_c$, 


$C_2=\sqrt{\frac{\mes\Omega}{m_0}}$, $\beta_j=DR^j$ ($j=1,2$), 


$\beta_3=\wt Q^2$, $\beta_4=\wt Q^1$. 


Тогда в силу неравенства (35), леммы 2 из (45) будем иметь 


$$ 


N(\varepsilon,t)\ge\tfrac{m_0}{2}\|\varepsilon\|^2_{L_2}- 


c[\|\varepsilon\|_{L_2}-\tfrac{1}{\delta} 


(\|\varepsilon_1\|_{L_2}+\|\varepsilon_2\|_{L_2}+1)]. 


$$ 


Отсюда следует, что при произвольно фиксированных 


$\|\varepsilon_j\|_{L_2}=r_j$ ($j=1,2$) радиус сферы 


$\|\varepsilon\|_{L_2}=r$ можем взять так, чтобы имело место 


\begin{equation} 


N(\varepsilon,t)\ge\frac{1}{3}\,\frac{m_0}{M_0}R^2,\quad 


0\le t\le 1. 


\end{equation} 


Из неравенств (42), (46) следует утверждение теоремы 4. 


 


Из теоремы 4 следует, что вполне непрерывное поле 


$\varepsilon-G(\varepsilon)$ гомотопно полю $\varepsilon$ 


на сферах достаточно большого радиуса пространства 


$\wt L_2(\ov\Omega)$, следовательно, 


его вращение на этих сферах равно $+1$. Тогда (\cite{1}, с.\,74) 


уравнение (31) внутри сферы $S_{\wt L_2}(R,0)$ 


имеет по крайней мере одно решение 


$\varepsilon=(\varepsilon_1,\varepsilon_2,\varepsilon_3)$, 


принадлежащее пространству $\wt L_2(\ov\Omega)$. 


 


 


\begin{thebibliography}{9} 


 


\bibitem{1} 


Ворович И.И. {\em Математические проблемы нелинейной теории пологих 


оболочек}. -- М.: Наука, 1989. -- 373~с. 


\bibitem{2} 


Тимергалиев С.Н. {\em Доказательство разрешимости одной задачи 


нелинейной теории пологих оболочек} // Изв. вузов. Математика. -- 1996. 


-- \No\,9. -- С.\,60--70. 


\bibitem{3} 


Векуа И.Н. {\em Обобщенные аналитические функции}. -- 2-е изд. -- М.: 


Наука, 1988. -- 509~с. 


 


 


\end{thebibliography} 


 


\vskip 2cm 


 


\post{\it Казанская государственная}{\it Поступила} 


\post{\it архитектурно-строительная}{27.04.1996} 


\post{\it академия}{} 


\smallskip 


\post{\it Камский политехнический институт}{} 


 


 


\label{end} 


\end{document} 


